% !TEX root = main.tex
\section{Number Theory Tools}\label{sec:number-theory-tools}
To analysis the relationship between the reduces and the items, we introduced some notations from number theory.

Since we used a lot of $a\mid b$ and $a\nmid b$ but they are not in $\mathbb{Z}$, at first, we have to clarify the definition of divisibility.
\begin{definition}\label{def:extend-divisibility}
Given two numbers $\frac{p_1}{q_1}, \frac{p_2}{q_2} \in \mathbb{Q}[X]$ and they are irreducible, define $\frac{p_1}{q_1}\mid \frac{p_2}{q_2}$ as $\exists t \in \mathbb{Z}_{\ne0}$ such that $t\cdot \frac{p_1}{q_1} = p_2$. If not, we say $\frac{p_1}{q_1}\nmid \frac{p_2}{q_2}$. Of course we stipulate that $\forall \frac{p_1}{q_1} \in \mathbb{Q}[X]$, $\frac{p_1}{q_1}\mid 0$. For $\mathbb{N}$, we treat as that $q = 1$.
\end{definition}
\begin{proposition}
    $\frac{p_1}{q_1}\mid \frac{p_2}{q_2} \Longleftrightarrow p_1\mid p_2$.
\end{proposition}
\begin{proposition}
    If $a \mid c$ and $b \mid c$, and $\gcd(a, b) = 1$, then $ ab \mid c$.
\end{proposition}
Below are number theory tools we introduced in this paper and their propositions.
\begin{definition}\label{def:order-under-modulo}
    Given two number $a, b\in \mathbb{N}_{_{>1}}$ that $\gcd(a, b) = 1$, denote
    \[\ord_{a}b = \min_{x \in \mathbb{N}_{_{\ge 1}}}\left\{b^{x} \equiv 1 \pmod{a} \right\}\]
    called the discrete order of $b$ modulo $a$.
\end{definition}
\begin{proposition}
    $\ord_{a}1 = 1$.
\end{proposition}
\begin{proposition}
    If $c \mid a$, then $\ord_{c}b \mid \ord_{a}b$.
\end{proposition}
\begin{proposition}\label{prop:power-order-mod}
    $\frac{b^{s\ord_{a} b} - 1}{b^{\ord_{a} b} - 1} \equiv s \pmod{a}$.
\end{proposition}
\begin{definition}\label{def:log-under-modulo}
Given $x\in \mathbb{N}_{>1}$.
If there exists $y \in \mathbb{N}_{_{\ge 1}}$ such that $a^y \equiv x \pmod{b}$ and $y \le \ord_{b}a$, then we denote $\log_{a, b}x = y$.
If there not exists a such $y$, we will always regard $\log_{a, b}x = \infty$.
\end{definition}
\begin{proposition}
    $a^{\log_{a, b}x} \equiv x \pmod{b}$.
\end{proposition}
\begin{proposition}
    $\forall n\in\mathbb{N}$, we have $\log_{a, b}(x + nb) = \log_{a,b}x$.
\end{proposition}
\begin{proposition}
    $\forall n \in \mathbb{N}$ we have $\log_{a+nb, b}x = \log_{a, b}x$.
\end{proposition}
\begin{proposition}
    Suppose $\log_{a, b}x < \infty$ and $\log_{a, b}y < \infty$, then we have
    \[ \log_{a, b}(x \cdot y) = (\log_{a, b}x + \log_{a, b}y) \pmod{\ord_{b}a}. \]
\end{proposition}
\begin{proposition}
    Suppose $\log_{a, b}x < \infty$.
    Given $n \in \mathbb{N}$, we have
    \[ \log_{a, b}(x ^{n}) = n\log_{a, b}x \pmod{\ord_{b}a}. \]
\end{proposition}
\begin{proposition}
    Suppose $\log_{a,b}x < \infty$, then $\log_{a, b}(x^{-1}) = \ord_{b}a - \log_{a,b}x$.
\end{proposition}
\begin{definition}\label{def:log-set-under-modulo}
If $\log_{a, b}x < \infty$, then we denote that
\begin{align*}
    \Log_{a, b}x &= \ord_b(a)\mathbb{Z} + \log_{a, b}x, \\
    - \Log_{a, b}x &= \ord_b(a)\mathbb{Z} - \log_{a, b}x.
\end{align*}
If $\log_{a, b}x = \infty$, we denote $\Log_{a, b}x = \varnothing$.
\end{definition}
We call $\log_{a, b}(x)$ the principal logarithm of $x$ on $b$ modulo $a$, and $\Log_{a, b}(x)$ the logarithms of $x$ based on $b$ modulo $a$, and we will define $\log_{a, b}1 = 0$ and $\Log_{a, b}1 = \ord_{b}(a)\mathbb{Z}$.
Without special state, we will always suppose the $x > 1$ in $\log_{a, b}x$ and $\Log_{a,b}x$.

\begin{proposition}
    $\forall n\in \Log_{a,b}x$, we have $a^{n} \equiv x \pmod{b}$.
\end{proposition}
\begin{proposition}
    $\Log_{a,b}(x^{-1}) = -\Log_{a,b}x$.
\end{proposition}
\begin{definition}\label{def:generated-subgroup}
    Define $(a;b) = \{ a^{x} \mid 0 \le x < \ord_{b}(a) \}$ as the generated subgroup of $a$ modulo $b$.
\end{definition}
\begin{definition}
    Given $A\subset \mathbb{Z}$ and $a\in \mathbb{N}$. We define $A/a = \frac{A}{a} = \{ x\fmod a\mid x\in A \}$ as the remainders of $A$ modulo $a$.
\end{definition}
\begin{proposition}
    $(a^{-1};b)/b = (a;b)/b$, and $(a;b)/b \le (\mathbb{Z}/b)^{*}$ by multiplication.
\end{proposition}
% \begin{proof}
%     $a^{x} \cdot a^{y} = a^{x + y} \equiv a^{(x + y) \fmod\ord_{b}a} \pmod{b}$ and $a^{(x + y) \fmod\ord_{b}a}\in (a;b)$. Besides, $a^{-x} \equiv a^{\ord_{b} a - x} \pmod{b}$ and $a^{\ord_{b} a - x}\in (a;b)/b$. 
% \end{proof}
\begin{definition}\label{def:power-and-power-count}Define $\val_{a}b = \max_{n \in \mathbb{N}}\{n: a^n \mid b\}$ as the exponent of maximum power of $a$ in $b$ on $a$.
    Define $\pow(a\mid b) = a^{\val_{a}b}$ as the maximum power of $a$ in $b$. 
    And we define $\rel(a\mid b) = \frac{b}{\pow(a\mid b)}$ called the remainder part of $b$ with respect to power of $a$.
\end{definition}
\begin{proposition}
    If $c \mid a$ and $a \mid b$, then $\val_{c}b \ge \val_{c}a \cdot \val_{a}b$.
\end{proposition}
\begin{definition}\label{def:radical}
    Define $\rad a = \prod_{p \in \mathbb{P}}^{p \mid a}p$ as the radical of $a$. Given $S\subset\mathbb{N}$, define $\Omega(S\mid a) = \{p \in S: p \mid a\}$ as all factors of $a$ in $S$. If $S = \mathcal{P}$, we call $\Omega(\mathbb{P}\mid a)$ radical set of $a$.
\end{definition}
\begin{definition}\label{def:filtered-radical-powers}
    Define $\raf_{a}b = \prod_{p \in \Omega(\mathbb{P}\mid a)}\pow(p\mid b)$ as the filtered radical power of $b$ with respect to $a$. And we also call it the $a$-filtered radical powers of $b$ or the radical $a$-part of $b$. Specially, we define $\raf_{1}b = 1$.
    And same like to the $\rel$,  we define $\ram_{a}b = \frac{b}{\raf_{a}b}$ called the remainder part of $b$ with respect to radical $a$-part.
\end{definition}
\begin{proposition}
    $\rad a = \prod_{p \in \Omega(\mathbb{P}\mid a)}p$
\end{proposition}
\begin{proposition}\label{prop:self-containment}
    $\raf_{a}b \mid b$.
\end{proposition}
\begin{proposition}
    $\raf_{a}a = a$.
\end{proposition}
\begin{proposition}\label{prop:raf-idempotency}
    $\raf_{a}\raf_{a}b = \raf_{a}b$.
\end{proposition}
\begin{proposition}\label{prop:raf-completely-multiplicative}
    $\raf_{a}(b_1b_2) = \raf_{a}b_1\cdot\raf_{a}b_2$.
\end{proposition}
\begin{proposition}\label{prop:raf-inclusion-exclusion-principle}
    $\raf_{a_1 a_2}b = \frac{\raf_{a_1}b \cdot \raf_{a_2}b}{\raf_{\gcd(a_1, a_2)}b}$.
\end{proposition}
\begin{proposition}\label{prop:raf-power-homogeneity}
    $\raf_{a}(b^n) = (\raf_{a}b)^{n}$.
\end{proposition}
\begin{proposition}\label{prop:raf-index-monotonicity}
    $a_1 \mid a_2 \implies \raf_{a_1}b \mid \raf_{a_2}b$.
\end{proposition}
\begin{proposition}\label{prop:raf-valuation-monotonicity}
    $\val_{p}b_1 \le \val_{p}b_2 \Longleftrightarrow \raf_{p}b_1 \mid \raf_{p}b_2$.
\end{proposition}
\begin{proposition}\label{prop:raf-radical-invariance}
    $\rad a_1 = \rad a_2 \implies \raf_{a_1}b = \raf_{a_2}b$.
\end{proposition}
\begin{proposition}\label{prop:raf-gcd-distributivity}
    $\raf_{a}\gcd(b_1, b_2) = \min\{\raf_{a}b_1, \raf_{a}b_2\}$.
\end{proposition}
\begin{proposition}\label{prop:raf-lcm-distributivity}
    $\raf_{a}\lcm(b_1, b_2) = \max\{\raf_{a}b_1, \raf_{a}b_2\}$.
\end{proposition}
\begin{proposition}\label{prop:raf-limit-power-gcd-raf-equvalent}
    $\lim_{n\to\infty}\gcd(a^{n}, b) = \raf_{a}b$.
\end{proposition}
\begin{proposition}
    $\log \raf_{a}b = \sum_{p\in\mathbb{P}}^{p \mid a} \val_{p}b \log p$
\end{proposition}
\begin{proposition}\label{prop:rem-coprimality}
    $\gcd(\ram_{a}b, a) = 1$.
\end{proposition}
\begin{proposition}\label{prop:raf-ord-non-self-divisibility}
    If $p \in \mathbb{P}$, $p\nmid b$, then $\raf_{p}\ord_{p}b = 1$.
\end{proposition}
\begin{theorem}\label{thm:lte-radical-theorem}
    If $p \in \mathbb{P}$ and $p\mid (b - 1)$, then
    \[\raf_{p}(b^{n} - 1) = \left\{ \begin{aligned}
        &\raf_{p}(b - 1) \cdot \raf_{p}n, & p > 2, \\
        &\raf_{2}(b - 1) \cdot \raf_{2}n + \raf_{2}(b + 1) - 1, & p = 2.
    \end{aligned} \right.\]
\end{theorem}

\begin{definition}\label{def:radicable-order}
    $\rord_{a} b = \lcm_{p\in\mathbb{P}}^{p\mid a} \{\ord_{p} b\}$ be the radicable order.
\end{definition}

\begin{definition}\label{def:radical-height-and-width}
    Given $a \in \mathbb{N}_{\ge 1}$. Denote $\nu(a) = \max\{ \val_{p} a \mid p\in\mathbb{P}\}$ and $\omega(a) = \card\Omega(\mathbb{P}\mid a)$ called the radical height and the radical width of $a$.
\end{definition}

% Now let's introduce a kind of count function on infinite subsets of $\mathbb{N}$ from analytic number theory. Given $A \subset N$ that $\card{A} = \infty$ and $\forall a\in A$, $a > 1$.
% \begin{definition}
%     Define $\pi(A;x) = \card\{ A \le x \}$ that $x\in\mathbb{R}_{>1}$. And we define the Perron form
%     \[\widetilde{\pi}(A;x) = \left\{\begin{aligned}
%     &\pi(A;x), && x\not\in A, \\
%     &\frac{1}{2}\left[\pi(A;x^{-}) + \pi(A;x^{+})\right], && x\in A.
% \end{aligned}\right.\]
% \end{definition}
% \begin{proposition}
%     $\sum_{t\in A}^{t\le x}f(t) = \int_{1}^{x}f(t)\rmd{\pi(A;t)}$.
% \end{proposition}
% \begin{proposition}
%     $\pi(A\cup B;x) = \pi(A;x) + \pi(B; x) - \pi(A\cap B; x)$. So does for $\widetilde{\pi}(A;x)$.
% \end{proposition}
% \begin{definition}
%     Define $D(A;s) = \sum_{x\in A}x^{-s}$ as the Dirichlet series of $A$.
% \end{definition}
% \begin{proposition}
%     $D(aA;s) = a^{-s}D(A;s)$.
% \end{proposition}
% \begin{proposition}
%     If $A\otimes B$ is uniqe decomposed, then we have
%     \[ D(A \otimes B; s) = D(A;s)\cdot D(B;s). \]
% \end{proposition}
% \begin{theorem}\label{thm:count-function-from-dirichlet-series}
%     For large enough $c > 0$, we have
%     \[ \widetilde{\pi}(A;x) = \frac{1}{2\pi\rmi}\int_{c-\rmi\infty}^{c+\rmi\infty}D(A;s)\frac{x^{s}}{s}\rmd{s}. \]
% \end{theorem}
% \begin{lemma}[Perron's formula]
% For $y > 0$, we have
% \[ P(x) = \frac{1}{2\pi\rmi}\int_{c-\rmi\infty}^{c+\rmi\infty}\frac{x^{s}}{s}\rmd s = \left\{\begin{aligned}
%     &1, &&y > 1,\\
%     &1/2, &&y = 1,\\
%     &0, &&y < 1.
% \end{aligned}\right. \]
% \end{lemma}
% \begin{proof}[Proof of \cref{thm:count-function-from-dirichlet-series}]
%     \begin{align*}
%         &\frac{1}{2\pi\rmi}\int_{c-\rmi\infty}^{c+\rmi\infty}D(A;s)\frac{x^{s}}{s}\rmd s \\
%         =&\frac{1}{2\pi\rmi}\int_{c-\rmi\infty}^{c+\rmi\infty}\sum_{t\in A}t^{-s}\frac{x^{s}}{s}\rmd s \\
%         =&\sum_{t\in A}\frac{1}{2\pi\rmi}\int_{c-\rmi\infty}^{c+\rmi\infty}\left(\frac{x}{t}\right)^{s}\frac{1}{s}\rmd s \\
%         =&\sum_{t\in A}P(x/t) \\
%         =&\widetilde{\pi}(A;x).
%     \end{align*}
% \end{proof}
% \begin{definition}
%     Define $\Pi(A;x) = \sum_{m=1}^{\infty}\frac{1}{m}\pi(A;x^{1/m})$ and also the Perron form $\widetilde{\Pi}(A;x) = \sum_{m=1}^{\infty}\frac{1}{m}\widetilde{\pi}(A;x^{1/m})$.
% \end{definition}
% \begin{definition}
%     For $s > 0$, we define $L(A;s) = \prod_{x\in A}(1-x^{-s})^{-1}$.
% \end{definition}
% \begin{proposition}
%     $\ln L(A;s) = \sum_{m=1}^{\infty}\frac{1}{m}\sum_{x\in A}x^{-sm} = s\int_{1}^{\infty}\frac{\Pi(A;x)}{x^{s+1}}\rmd x$.
% \end{proposition}
% \begin{proof}
% At first, we have
% \[ \ln L(A;s) = -\sum_{x\in A}\ln(1 - x^{-s}) = \sum_{x\in A}-\ln(1 - x^{-s}). \]
% From Maclaurin's formula, $- \ln (1 - z) = \sum_{m=1}^{\infty}\frac{x^{m}}{m}$, we have
% \begin{align*}
%     \ln L(A;s) &= \sum_{x\in A}\sum_{m=1}^{\infty}\frac{x^{-sm}}{m} = \sum_{m=1}^{\infty}\frac{1}{m}\sum_{x\in A}x^{-sm} \\
%     &= \sum_{m=1}^{\infty}\frac{1}{m}\int_{1}^{\infty}x^{-sm}\rmd\pi(A;x) \\
%     &= \sum_{m=1}^{\infty}\frac{1}{m}\int_{1}^{\infty}x^{-s}\rmd\pi(A;x^{1/m}) \\
%     &= \int_{1}^{\infty}x^{-s}\rmd\left(\sum_{m=1}^{\infty}\frac{1}{m}\pi(A;x^{1/m})\right) \\
%     &= \int_{1}^{\infty}x^{-s}\rmd\Pi(A;x)
% \end{align*}
% According the rule of integration by parts, we have
% \begin{align*}
%     \ln L(A;s) &= \left[x^{-s}\Pi(A;x)\right]\bigg\vert_{1}^{\infty} - \int_{1}^{\infty}\Pi(A;x)\rmd x^{-s} = s\int_{1}^{\infty}\frac{\Pi(A;x)}{x^{s+1}}\rmd x.
% \end{align*}
% \end{proof}
% \begin{theorem}
%     Let $\mu(m)$ be the Möbius function, for $c > 0$ large enough, we have
%     \[ \widetilde{\pi}(A;x) = \frac{1}{2\pi\rmi}\sum_{m=1}^{\infty}\frac{\mu(m)}{m}\int_{c-\rmi\infty}^{c+\rmi\infty}\ln{L(A;ms)}\frac{x^{s}}{s}\rmd s. \]
% \end{theorem}
% \begin{proof}
%     Since we have Perron's formula, we are supported to have
% \begin{align*}
%     &\frac{1}{2\pi\rmi}\int_{c-\rmi\infty}^{c+\rmi\infty}\ln{L(A;s)}\frac{x^{s}}{s}\rmd s \\
%     =&\frac{1}{2\pi\rmi}\int_{c-\rmi\infty}^{c+\rmi\infty}\sum_{t\in A}\sum_{m=1}^{\infty}\frac{t^{-sm}}{m}\left(\frac{x^{s}}{s}\right)\rmd s \\
%     =&\sum_{m=1}^{\infty}\sum_{t\in A}\frac{1}{m}\left[ \frac{1}{2\pi\rmi}\int_{c-\rmi\infty}^{c+\rmi\infty}\left(\frac{x}{t^{m}}\right)^{s}\frac{1}{s}\rmd s \right] \\
%     =&\sum_{m=1}^{\infty}\sum_{t\in A}\frac{1}{m}P(x/t^{m}) \\
%     % =&\sum_{m=1}^{\infty}\frac{1}{m}\card\{t\in A \mid \frac{x}{t^{m}} > 1 \} \\
%     % =&\sum_{m=1}^{\infty}\frac{1}{m}\card\{t\in A \mid t^{m} < x \} \\
%     =&\sum_{m=1}^{\infty}\frac{1}{m}\widetilde{\pi}(A;x^{1/m}) \\
%     =&\widetilde{\Pi}(A;x).
% \end{align*}
% Using the Möbius inversion, we have
% \begin{align*}
%     \widetilde{\pi}(A;x) &= \sum_{m=1}^{\infty}\frac{\mu(m)}{m}\widetilde{\Pi}(A;x^{1/m}) \\
%     &= \sum_{m=1}^{\infty}\frac{\mu(m)}{m}\left[ \frac{1}{2\pi\rmi}\int_{c-\rmi\infty}^{c+\rmi\infty}\ln{L(A;s)}\frac{x^{s/m}}{s}\rmd s \right] \\
%     &= \frac{1}{2\pi\rmi}\sum_{m=1}^{\infty}\frac{\mu(m)}{m}\int_{c-\rmi\infty}^{c+\rmi\infty}\ln{L(A;ms)}\frac{x^{s}}{s}\rmd s.
% \end{align*}
% \end{proof}